* Recaptulação: lista as contas do PL
** Patrimômio Líquido
*** Capital Social
****  Capital Social Integralizado
**** (-) Capital social a integralizar
*** Reserva de Capital
**** Ágio na venda de ações
**** (-) custos de distribuição
*** Reserva de lucro
**** Reserva legal
**** Reservas de Inscentivos Fiscais
**** Reservas Estatutárias
**** Reservas de Contingência
*** Ajustes de Avaliação patriminial
*** (-) Ações em tesouraria
*** Prejuizos Acumulados
*** Reserva de Lucro (não tem em SA)
*** Dividendos adicionais propostos.
* Exercicio



| Conta               | Exercicio 0 | Exercicio 1 |
|---------------------+-------------+-------------|
| Capital Social      |      800000 |             |
| Reserva Legal       |      120000 |     145.000 |
| REserva Eepansão    |      200000 |     300.000 |
| Reserva Estatutária |      100000 |       50.00 |
| Divi. Adic. Propo.  |             |     182.500 |
|---------------------+-------------+-------------|
| Total               |     1120000 |             |
#+TBLFM: @>$2=vsum(@2..@-1) 



| Item                   | composição     | Referencia 1 | Referencia 2 |   Valor |
| Lucro liquido          |                |              |              | 500000. |
| Reserva Já constituida |                |              |              | 120000. |
| LImite I               | 500000 * 5%    |           5% |              |  25000. |
| LIMITE ii              | 800000 * 30%   |       800000 |          30% | 240000. |
|                        | Limite II - R. |              |              | 120000. |
| reserva constituida    | 120000         |       228000 |              |         |
#+TBLFM: $5=@-1-@-3::@4$5=@-2*$-2



| movimentação  | Cap. Social | R. Legal | R. Estat. | Res. Ex. | ORA | Div. Obrig. | Total |
| Saldo Inicial |             |          |           |          |     |             |       |
|               |             |          |           |          |     |             |       |
* 
* 
